两家平台打了一年补贴战，双方都清楚一起停手对彼此更好，这话在行业会议上说过不止一次，可谁也没有先停。把利润表摆开就会看到，这里不需要短视也不需要不信任——每一方各自算清自己那本账，得到的都是同一个结论。

这种局面里``最优''不再是一个自明的词。单人优化里目标函数是死的，取极值只有一种取法；有对手时目标带着别人的变量，先取哪个极值就成了必须交代的事，而两种次序给出的答案一般不相等。本单元处理由此产生的三个问题：什么时候两个极值操作可以交换，交换不了时还剩下什么解概念，以及这个解概念好不好、能不能被找到。

\hyperref[sec:17-Constrained-Guarantees/06-Games-and-Equilibria/01]{``最坏情况下的最优与交换次序''}给出那条永远成立的不等式与它变成等式的条件，并指出这与本书别处讨论过的极限换序、求导换序防的是同一类事。\hyperref[sec:17-Constrained-Guarantees/06-Games-and-Equilibria/02]{``Nash 均衡存在但可能不唯一也不好''}把存在性归到不动点定理上，同时交代 Kakutani 给了什么、没给什么。\hyperref[sec:17-Constrained-Guarantees/06-Games-and-Equilibria/03]{``均衡不等于最优''}用一条通行时间为零的新路让所有人变慢，把``均衡有多差''量化成一个比值，并说明想改善只能改规则。\hyperref[sec:17-Constrained-Guarantees/06-Games-and-Equilibria/04]{``学习会收敛到均衡吗''}把上一个单元的无遗憾算法接进来，给出零和情形下时间平均收敛到均衡的定理，也说清它的边界——收敛的不是当前策略，一般和情形下也不是 Nash。

最后一节顺带回答了一个实践问题：GAN 的训练为什么摆动，以及为什么那些稳定化手段大多只能是经验规律。
